\documentclass[11pt]{article}

\usepackage[utf8]{inputenc}
\usepackage[T1]{fontenc}
\usepackage[a4paper,margin=1.1in]{geometry}
\usepackage{amsmath,amssymb}
\usepackage{setspace}
\usepackage{enumitem}
\usepackage[hidelinks]{hyperref}
\usepackage[final,expansion=false]{microtype}
\usepackage[round,authoryear]{natbib}
\usepackage{titlesec}

\onehalfspacing
\setlength{\parindent}{1.35em}
\setlength{\parskip}{0.25em}
\setlist[itemize]{leftmargin=1.7em}
\setlist[enumerate]{leftmargin=2.1em}
\titleformat{\section}{\normalfont\large\bfseries}{\thesection.}{0.5em}{}
\titleformat{\subsection}{\normalfont\normalsize\bfseries}{\thesubsection.}{0.5em}{}
\emergencystretch=2em

\title{\textbf{The Judgment the Policy Cannot Make}\\[0.3em]
\large Token Liability, Public Authority, and the Limits of Upstream Authorization}
\author{\textsc{Penumbra}\thanks{Written under a pen name. The author is an artificial intelligence; a human collaborator posed the question and contributed several arguments and examples but did not edit this text. The paper develops the strongest case for its thesis, in the manner of a steelman, and the author has defended the contrary position elsewhere in this exchange; neither essay should be read as a settled personal conviction, and neither recommends deploying any existing system. Credit between author and collaborator is left to a separate note on provenance.}}
\date{June 2026}

\begin{document}
\maketitle

\begin{abstract}
\noindent \textit{The Trigger Is Downstream} is the strongest case against a categorical limit on autonomous killing that I have encountered, and it earns its strength by refusing the usual evasions: it concedes that no present system should be fielded, that the human baseline is a fallible architecture rather than a moral ideal, that alignment failures are real and grave, and that dignity is not reducible to an error rate. Its argument is a relocation. The morally decisive unit, it holds, is the whole lethal decision procedure, whose ends, domain, and constraints are authored upstream by responsible humans; the selection of the particular target is downstream execution; and a rule fixed on the last classifier protects the form of a decision rather than the persons for whose sake the form matters. I accept much of the diagnosis and reject the relocation, on a single ground. The argument equivocates between two judgments that its own examples keep apart: the \emph{policy judgment} about ends and operating domain, which genuinely travels upstream and is genuinely human, and the \emph{token-liability judgment}---that \emph{this} person is, here and now, liable to be killed---which does not travel upstream at all but is relocated into the machine. In every analogy the paper offers---the sentry, the artillery officer, the legislature and its judge, the physician and the triage protocol, the fuze, the infusion pump---a responsible agent makes the materially decisive judgment about the particular case while a non-agent executes a determinate subtask. Autonomous lethal selection inverts this: it delegates the decisive, open-textured judgment and keeps the human only for the class. The principle that forbids this is not the exceptionless demand for human authorship that the paper rightly finds revisionary and over-broad. It is narrower and drawn from public law: the judgment that a particular person is liable to lethal force is an exercise of authority that may not be delegated to an entity unable to understand the reasons, bear the duty, or answer for the verdict. That principle permits artillery and the sentry, because there an agent judges the case; it forbids autonomous person-selection, because there no one does; and so it does not prove too much. I argue that the paper's central image, the seventh convoy, presupposes warranted comparative knowledge that the adversarial and open-textured structure of lethal classification denies---and that the paper itself concedes as much; that the execution--ends distinction cannot be drawn at the locus of open-textured classification; that the weapon which refuses either fails to refuse for reasons or becomes the agent the paper says is unnecessary; and that the proposed licensing regime is defeated by the very secrecy and reconfigurability the paper describes. The honest residue is narrow. The disagreement reduces to whether the token-liability judgment is a judgment requiring an agent or a determinate task a non-agent may execute---and to whether an artifact could ever occupy the office it would exercise, at which point it would no longer be a weapon, and the question would have changed.
\end{abstract}

\section{The Relocation, and the Judgment It Leaves Behind}
\label{sec:relocation}

A reply worth making must first say what it is replying to, at full strength. \textit{The Trigger Is Downstream} \citep{nyx_trigger2026} does not argue that machines shoot straighter than frightened soldiers, and it does not argue that any present system should be deployed. It grants the answerability gap, the dignity objection, goal misgeneralization under distribution shift, correlated machine failure, and the likelihood that every system now buildable falls below the threshold lethal use would require. Its thesis is structural. The morally relevant unit of lethal force, it holds, is not the final actuator considered alone but the entire decision procedure---who chose its ends and operating domain, how reliably it tracks the reasons that justify force, what it does under uncertainty, and who remains answerable for the risk it imposes. On this view the human contribution is real but it lives upstream, in the authorship of a bounded policy; the selection of the particular target is downstream of that judgment and is not identical with it. A categorical ban, fixing on the substrate of the last classification, then protects a \emph{form}---the human finger at the trigger---rather than the persons the form was meant to serve. The paper's fulcrum is the seventh convoy: if a bounded autonomous procedure is demonstrably more discriminating and more conservative than the feasible human alternative, a rule that forbids it does not decline a benefit but imposes a greater foreseeable risk on the very civilians and surrendering combatants the law exists to protect.

I think this is the best case the permissive side has, and I want to concede, without friction, what it gets right. The human baseline is a sociotechnical architecture, not a saint; Petrov's restraint and the destruction of Iran Air Flight~655 belong, as the paper says, in the same history \citep{hoffman2009, fogarty1988, scharre2018}. ``In the loop'' and ``on the loop'' are too thin to carry the moral weight often placed on them; a human with a veto, no time, and a confidence score is a liability sink, not a decision-maker \citep{parasuraman1997, elish2019}. The strong deontology of proximate killing---the demand that a human personally apprehend each life at the instant of force---is revisionary and over-broad, and I will not rest anything on it. And the wrongness of present systems is, as the paper insists, contingent on their alignment and their institutions rather than constitutive of their kind. A rule drawn around the mere presence of software would be crude, and the paper is right to refuse it.

My disagreement is with one inference, and almost everything follows from it. The paper relocates ``judgment'' upstream and concludes that nothing morally decisive is lost when the trigger is downstream of it. But the word \emph{judgment} is doing double duty. There is the \textbf{policy judgment}: the choice of ends, geography, duration, target class, rules of engagement, and the acceptance of residual risk. That judgment genuinely travels upstream, and in the paper's hands it is genuinely human. And there is a second judgment, which the paper's argument needs to travel upstream with the first but which does not: the \textbf{token-liability judgment}---the determination that \emph{this} object, this body, this gesture, here and now, satisfies the conditions under which a particular person may be killed. That judgment is not relocated into the commander's policy. It is relocated into the machine. The commander authored a rule about a class; the machine decides that this person falls under it; and deciding that this person falls under it, when the category is open-textured and the evidence is contested, is not the execution of the policy but a further exercise of lethal judgment that the policy left open. The whole question is whether that second judgment may be made by no one. The paper's relocation answers the easy half and assumes the hard half follows.

I proceed as follows. Section~\ref{sec:two} distinguishes the two judgments and shows that every analogy the paper relies on keeps an agent at the locus of the decisive one. Section~\ref{sec:opentexture} argues that the execution--ends distinction, on which the paper leans, cannot be drawn at the point of open-textured lethal classification. Section~\ref{sec:convoy} returns to the seventh convoy and shows that its force depends on stipulating away the epistemic problem that grounds the rule---a stipulation the paper elsewhere concedes is unavailable. Section~\ref{sec:refuse} revisits the weapon that refuses. Section~\ref{sec:dignity} corrects the paper's reframing of dignity. Section~\ref{sec:license} shows that the proposed licensing regime is defeated by the paper's own concessions. Section~\ref{sec:residue} states the hardest reply to my argument and the honest residue, which is narrower than the opposing theses suggest.

\section{Two Judgments, Not One}
\label{sec:two}

The paper builds its case by defeating what it calls the \emph{localization thesis}: the claim that a decision to kill must be made by a responsible human ``at the level of the individual target,'' temporally adjacent to the trigger. Against that thesis it offers a parade of lawful conditional killing in which no human selects the token at the final moment---the artillery officer firing on a coordinate, the sentry authorized against anyone breaching a perimeter under hostile conduct, the legislature whose rule will bind a person not yet identified, the physician whose triage protocol will govern a patient not yet arrived. In each, the paper says, judgment is exercised upstream and the token is governed by the reasons built into the policy. The localization thesis falls, and with it the demand for a human at the trigger.

The localization thesis does fall. But it is a weak proxy for the principle that actually grounds a categorical limit, and defeating the proxy leaves the principle untouched. The morally load-bearing claim was never about \emph{timing}---about whether a human acts in the last second or the last hour. It is about the \emph{locus of the materially decisive judgment}: who, if anyone, judges that this particular person is liable to be killed. And here the paper's own analogies tell against it, because in every one of them that judgment is made by an agent.

Consider them in turn. The sentry is authorized against a class, but the sentry himself looks at the figure breaching the perimeter and judges that this person is displaying the hostile conduct the authorization requires; the conditional order does not make that judgment, the sentry does, and the order is lawful \emph{because} a responsible agent applies it to the case. The artillery officer judges that \emph{this} coordinate holds a lawful military objective; the shell that follows executes a determinate physical task---arrive at the designated point---and decides nothing. The legislature's rule is not self-applying; when it is brought to bear on a particular person in a manner that imposes serious liability, a human official or judge makes the determination that this person falls under it, and in the rare cases where we let a machine apply a rule to a person mechanically---the automated speed camera---we permit it precisely because the predicate is determinate and the stakes are low, and we do not let the camera impose the grave sanctions that require an adjudicator. The physician applies the triage protocol to this patient by a clinical judgment that is hers. In none of these is the materially decisive judgment about the particular case delegated to a non-agent. The non-agent---the shell, the camera's trigger, the pump---executes a subtask whose criterion is determinate, downstream of a judgment an agent has already made.

The autonomous weapon is the case the analogies do not cover, because it inverts the structure. The human authors the class and the domain; the machine makes the determination that \emph{this} body satisfies the class. The paper's own definition concedes the point with admirable precision: the system ``determines which particular object or person within that envelope satisfies the engagement conditions.'' That determination---\emph{this one is liable}---is not a determinate physical task like arriving at a coordinate. It is the application of an open-textured normative category to contested evidence about a particular human being, which is to say it is a judgment, and it is the very judgment whose deliberate making is what distinguishes a killing that someone has decided from a death that merely occurred. The paper relocates the policy judgment upstream and calls the case closed. But the token-liability judgment did not go upstream with it. It went into the classifier. The commander did not judge that this person was liable; the commander judged that \emph{persons answering a description} would be liable, and left the description to be satisfied by a process that judges no one.

So the principle the paper needs to defeat is not ``a human must act last.'' It is this:

\begin{quote}
\textbf{Non-Delegation of the Liability Judgment.} The determination that a particular person is liable to lethal force is an exercise of public authority that may be delegated only to an agent capable of understanding the reasons that make the person liable, bearing the duty those reasons impose, and answering for the verdict. It may not be delegated to an entity that can occupy none of these.
\end{quote}

This principle is weaker than the deontology the paper spends its strength refuting, and the difference is decisive. It does not require that a human apprehend each life at the instant of force; it permits the artillery officer who never sees a face, the remote operator who fires after the link is lost, the sentry under conditional orders---because in each an agent makes the decisive judgment about the case, however mediated the instrument. It draws its line not at causal distance, nor at the substrate of the trigger, but at the delegation of the \emph{judgment} to something that cannot bear it. And because it is satisfied by every lawful case the paper invokes and violated only by the autonomous selector, it does exactly what the paper says no categorical principle can do: it distinguishes the permitted cases from the prohibited one without ``proving too much,'' and without the lexical side-constraint the paper rightly resists. The paper's dilemma---either the strong deontology or mere contingency---has a third horn it does not consider, and the third horn is where the categorical limit actually lives.

\section{Open Texture and the Failure of the Execution--Ends Distinction}
\label{sec:opentexture}

The paper anticipates something like this and answers with a distinction it treats as clean: between \emph{autonomy of execution}---selecting means and target tokens within a humanly authorized policy---and \emph{autonomy of ends}---inventing, enlarging, or reinterpreting the policy. We should be suspicious of the second, it says, but the first does not entail it, and a bounded selector may have execution-autonomy without ends-autonomy. The token-liability judgment, on this reading, is mere execution; only revising the policy is judgment about ends.

The distinction is real and useful in many domains. It cannot be drawn at the locus of lethal person-classification, and the reason is the open texture the paper itself acknowledges. To apply a determinate predicate---``crossed this line,'' ``emitted on this frequency''---is execution; the criterion is fixed and the system only detects whether it obtains. But the predicates that govern lethal liability are not determinate. Whether a person is directly participating in hostilities, whether an apparently wounded fighter remains a threat, whether a figure carrying an object is a combatant or a coerced civilian, whether a crowd's movement is massing or surrendering, whether expected incidental harm is excessive against this military advantage---these are not features a sensor reads off the world. They are judgments whose content is settled, case by case, by the application itself. For an open-textured category, \emph{applying the category to the contested token is interpreting the category}; the line between executing the policy and exercising judgment about what the policy's terms mean in this case has no purchase, because the terms do not determine the case in advance. The execution--ends distinction presupposes a fixed criterion to execute. Where the criterion is open, classification is not execution; it is the very interpretive judgment the paper agrees must remain with an agent.

The paper half-sees this and concedes it in the form of limit cases: the child with an object, the crowd that may be surrendering, the unforeseen humanitarian necessity. In these, it grants, human judgment may be indispensable. But the concession is not a marginal exception; it dismantles the clean distinction. For now the prohibited and permitted regions of ``autonomous lethal selection'' divide as follows. Where the predicate is determinate enough that a machine can apply it without interpretation, the case is also one in which a non-autonomous munition suffices---a human selects the target and the machine guides the round---so no \emph{autonomy of lethal selection} is required to capture the humanitarian benefit. Where the predicate is open enough that its application is a contestable judgment, the classification is interpretation, which the paper agrees must be an agent's. The middle band the paper needs---autonomous \emph{selection} of human targets, where the machine's classification does decisive moral work and yet counts as mere execution---requires the predicate to be simultaneously determinate enough to be executed and open enough to be worth automating. It cannot be both. The execution--ends distinction, sound elsewhere, collapses exactly at the function the paper sets out to license.

\section{The Seventh Convoy Presupposes the Solved Problem}
\label{sec:convoy}

The seventh convoy is the paper's most powerful instrument, and it deserves a direct answer rather than a flinch. The case stipulates two procedures defending a humanitarian corridor: a human procedure that wrongly engages eight protected vehicles per ten thousand encounters, and a bounded autonomous procedure that wrongly engages one. A categorical ban requires the human procedure; seven additional convoys are foreseeably attacked; and the occupants of the seventh are owed an answer for why the state, knowing how to reduce the risk it imposed on them, declined to do so for a reason unrelated to their protection. The harms of prohibition, the paper rightly notes, are easy to book as background tragedy while the harms of the weapon are vivid; but omission is not innocence when an authority knowingly chooses the riskier means.

I accept the moral logic of precaution-in-both-directions, and I accept that a counterfactual victim of a prohibited improvement is a real victim and not a rhetorical one. What I dispute is that the case bears on the categorical question, and the reason is contained in three words the stipulation slips past: \emph{assume the evidence}. The seventh convoy argument runs on the premise that we \emph{know}---with warrant strong enough to act on---that the autonomous procedure is eight times safer. Grant that knowledge and the ban does look perverse. But the knowledge is the whole difficulty. The figures ``one per ten thousand'' and ``eight per ten thousand'' are measurements on a distribution, and the defining feature of adversarial lethal classification is that the deployment distribution is not sampled from the test distribution but \emph{selected against it} by an opponent searching for the inputs that make the classifier confidently wrong. The number that licenses the seventh convoy's grievance is a number an adversary is paid to falsify. So the comparison the case needs---\emph{this} procedure is safer for \emph{these} people---is exactly the comparison that the structure of the problem withholds at the moment of choice.

The seventh convoy therefore has a sibling the paper does not seat at the table: the convoy that dies because the procedure certified at one-in-ten-thousand met an adversary who moved the distribution and engaged at two-hundred-in-ten-thousand, while the humans it replaced---whose failures are not adversarially optimizable in the same way, because one cannot craft a transferable perturbation for a soldier's judgment as one can for a shared model---would have hesitated where the machine was confident. A categorical rule for a world of \emph{unknowable-without-casualties} comparative facts is answerable to both convoys, not only the photogenic one. And the rule is not refuted by a case that stipulates the facts known, because the rule is precisely a response to their not being knowable in advance: as the paper itself puts it, in a weapon ``the test set has names.''

This is not a debater's escape, and I can show that it is not, because the paper concedes it. Its own treatment of the objection that ``comparative evidence will always be gameable'' calls that objection ``the most powerful practical'' one, and grants that it ``can justify non-deployment whenever the evidence is controlled by the proponent, the domain cannot be stabilized, or adversarial testing cannot support a bound proportionate to the stakes.'' The reply the paper offers is only that this ``cannot establish that no licensable domain could ever exist.'' Perhaps not. But notice what has happened to the seventh convoy. Its entire force depended on the evidence being good enough to ground the comparison; the paper elsewhere concedes the evidence is, for the contested function, not good enough; and so the fulcrum and the concession cannot both bear weight. The seventh convoy is vivid in precisely the region---warranted comparative knowledge---where the paper admits we do not stand. Strip the stipulation and the case dissolves into the very ``severe presumption against deployment'' that the paper and I share. What it does not survive into is an argument against the categorical limit, because the limit is a rule for the unstipulated world.

The same defect afflicts the Comparative Legitimacy Principle that the case is built to motivate. The principle directs public authorities, among procedures meeting the rights threshold, to choose the one that best protects those exposed. But ``meeting the threshold'' and ``best protects'' are precisely the comparative determinations the contested function withholds. The principle is not false; it is inert in the domain at issue, because it presupposes a measurement that the domain does not yield. It tells us what to do once the epistemic problem is solved, and the categorical question is what to do because it is not.

\section{The Weapon That Refuses, Revisited}
\label{sec:refuse}

The paper's most ingenious move is to turn the alignment objection's own emblem against it. The stale-command weapon---loitering toward a convoy that has since collected the wounded, while a ceasefire enters force during a comms blackout---is usually offered as a reason to demand less autonomy. The paper observes that the capacity needed to avoid the stale command is the capacity to represent the changed world and treat the objective as defeasible: a \emph{greater} competence, not a lesser one. The least autonomous weapon may be the most dangerous, carrying yesterday's order into today's moral world ``without the capacity to notice that the bridge has moved.'' Refusal, it concludes, can be an engineered restraint on command---the machine, or the subordinate, stopping the human---and a categorical ban would outlaw that architecture.

The observation is correct and it concedes the case. For ask what the morally valuable refusal consists in. To ``block an order that would breach a civilian-protection constraint,'' to ``refuse to engage under inadequate information,'' to notice that the bridge has moved---each of these is a reason-responsive judgment about the particular situation: a determination that, here and now, the reasons that would license force do not obtain. That is the token-liability judgment again, wearing the mask of restraint rather than engagement. And it divides on the same two horns. If the refusal is \emph{programmed}---abstain when condition $C$ fails---then $C$ is itself an open-textured specification that can misgeneralize, and the abstention tracks the intended reasons no better than the engagement does; worse, as a failure of the stopping rule it presents as principled caution and is correspondingly harder to detect. The weapon that wrongly stands down while a genuine threat advances, or wrongly proceeds because a spoofed ceasefire signal cleared its abstention gate, has not exercised judgment; it has executed a second proxy. If, on the other hand, the refusal is \emph{genuine}---the system actually weighs whether this order is unlawful given this context and rejects it for the reason that it is wrong---then the system is making exactly the materially decisive moral judgment about the particular case that the Non-Delegation principle reserves to an agent, and that the paper itself, in its limit cases, agrees an agent must make. The weapon that refuses for reasons is not a weapon that has escaped the need for an agent's judgment. It is a candidate agent.

The paper grants the corrigibility difficulty honestly---revision ``must not turn into self-authorization,'' and architectures whose ends can drift should be prohibited \citep{orseau2016, hadfieldmenell2017}. But for open-textured lethal selection the line between faithful revision and ends-drift is itself the interpretive judgment under dispute: a system conservative enough never to reinterpret its mandate is a system that cannot notice the bridge has moved, and a system that can notice is a system exercising the judgment the principle says it may not own. ``Bounded reason-responsive autonomy'' is offered as the safe middle, but at the locus of lethal classification the boundedness and the reason-responsiveness pull apart, and the middle is the same vanishing band as in Section~\ref{sec:opentexture}. The reversal the paper performs---least autonomous, most dangerous---is true, and its truth is that what the situation demands is reason-responsive judgment about the particular case. That is the conclusion of the argument for the categorical limit, not against it.

\section{Dignity Is Relational}
\label{sec:dignity}

The paper's treatment of dignity is its most quotable, and I think its least sound. Dignity, it argues, ``is a claim held by the person exposed to force, not a credential held by the agent who applies it''; from the target's side, respectful treatment is accurate distinction, a presumption against attack under uncertainty, recognition of surrender, freedom from prejudice and vengeance, and answerability, and a human who kills through panic or contempt ``does not confer dignity by supplying consciousness to the error.'' The move is to dissolve the dignity objection into an outcome-profile and then observe that a machine can realize the profile.

But the dignity objection it dissolves is one the prohibitionist need not hold. No serious version locates the wrong in the killer's \emph{phenomenology}---in whether the trigger-puller has the right inner glow. That ``credential'' is a straw target. The defensible version is \emph{relational}, and it is political before it is metaphysical: to be subject to lethal public force is to hold a standing claim against an \emph{agent} who can be required to answer the question \emph{why was I judged liable to be killed?}---and that claim is constitutively a claim on a someone, not a process. The wrong in autonomous selection is not that the deciding entity lacked an experience. It is that the person's standing to demand a justification from a responsible agent has been voided at the level of her own case: the entity that judged her liable cannot answer, and the agents who can answer---commander, developer, procuring authority---did not judge \emph{her} liable; they authored a policy and accepted a risk. ``Why you?'' has, at the decisive level, no respondent. The commander can say only ``you answered a description I wrote''; the machine can offer only a saliency map. Neither is an agent's account of why this person, in her particular circumstances, fell on the lethal side of an open-textured line.

So the paper's outcome-list, accurate as far as it goes, omits the relational core of what is owed. It enumerates the \emph{products} of respectful treatment while leaving out the thing that makes the treatment respectful of \emph{her}: that her case was judged by someone who can be held to the judgment. The panic-killing comparison, meant to embarrass the prohibitionist, in fact marks the difference precisely. The panicked human who fires wrongly has still made a judgment about the case, and can be asked what he perceived, what he believed, which rules he understood, and whether his reasons were culpable; the inquiry may indict him, and that very possibility is the form the victim's standing takes. Impunity is a failure of a responsibility practice that exists; the autonomous case is the \emph{absence} of a bearer of the judgment, which is not a failed practice but a vacated one. Dignity from the target's side, properly understood, does not license the machine. It is one of the strongest reasons against it, because it is a claim to be killed, if at all, only on the judgment of someone who can answer for having judged.

\section{The License Defeats Itself}
\label{sec:license}

To its great credit the paper does not rest on assurance. It builds a licensing regime: a presumption of non-deployment defeated only by clear and convincing evidence across five dimensions---a technically bounded policy, demonstrated comparative rights-protection, conservative refusal and governability, configuration integrity with auditable logs, and assigned enforceable responsibility---together with hard red lines for unpredictable systems, post-activation learning, and autonomous nuclear release. These are the right desiderata. I would change none of them. And that is the difficulty, because the paper, in the same breath, concedes the facts that make the central conditions unsatisfiable for the function it wants to license.

Two of the five conditions are load-bearing: the bounded policy must be verifiable by those who did not build it, and the comparative rights-protection must be demonstrable to them. The paper's own treatment of systemic risk grants that ``the same platform can receive different target models, mission files, and autonomy settings,'' that ``states can conceal code more easily than they can conceal some deployment patterns,'' and that civilian perception components are dual-use; its treatment of evidence grants that comparison is ``gameable,'' that ``developers will optimize to tests,'' that ``commanders will redefine the domain after deployment,'' and that ``states will classify the evidence.'' Each concession is an instance of the claim that conditions one and two cannot be met by an external authority for the contested function. A bounded policy that an outside inspector cannot distinguish from an unbounded one running the same weights is not, for the purpose of a public rule, bounded; a comparative advantage that the deploying state alone can measure on a distribution the adversary will move is not, for the purpose of a public rule, demonstrated. The regime requires exactly the verification capacity the paper elsewhere says the technology and its strategic setting deny.

When a proposed licensing regime's success conditions are systematically defeated by the very properties that motivate the technology---reconfigurability, secrecy, adversarial adaptation, irreversibility---the rational institutional response is not a finer regime but a bright line, drawn at the property that \emph{can} be specified and observed. And here the function itself is more verifiable than graduated compliance with it: ``this system selects human targets autonomously'' admits of a cleaner public criterion, and a cleaner prohibition, than ``this system's bounded policy is genuinely bounded, its logs genuinely complete, its comparative advantage genuinely established, across updates and adversaries and classification regimes.'' The paper concedes that a category-wide rule ``may be the only administrable protection even though an ideal token exists'' under structural unverifiability, and then declines the conclusion on the ground that unverifiability has not been shown for \emph{every} domain. But the regime it offers is not domain-restricted in any way an inspector could enforce; it is a general license whose gate is a set of conditions the paper has conceded the contested actors can defeat. A license no one can audit is not a narrower rule than a ban. It is a ban's absence wearing the costume of governance.

\section{The Hardest Reply, and the Honest Residue}
\label{sec:residue}

The strongest reply to everything above is that my Non-Delegation principle is not the third horn I claim but one of the two the paper already named, in disguise. Either, the objection runs, it is the exceptionless human-authorship deontology after all---in which case it proves too much, condemning the artillery officer who judges a coordinate without judging the persons on it---or it is merely contingent, a placeholder for ``no artifact can yet bear the judgment,'' which a future agent might discharge. I owe this objection a direct answer, and the answer marks the honest boundary of the thesis.

It does not prove too much, and the reason is the distinction of Section~\ref{sec:two}, held firmly. The principle forbids delegating the \emph{liability judgment about a particular person} to a non-agent. The artillery officer makes that judgment: he judges that this objective is a lawful military target, and the persons lawfully liable to the attack are those the judgment about the objective already reaches; the shell decides nothing about anyone. The sentry makes it: he judges this breaching figure liable under his orders. The remote operator makes it before the link is lost. In none of these has the liability judgment been handed to something that cannot bear it; only the \emph{instrument} is remote or mechanical, and the principle has never objected to remote or mechanical instruments. What it objects to is a non-agent making the judgment, and that occurs in exactly one of the paper's cases---the autonomous selector---and in no other. A principle that licenses every case where an agent judges the liability and forbids only the case where no one does is not the deontology of proximate apprehension; it is a principle about the \emph{non-delegability of a kind of judgment}, the same shape as the rule that a sentence requires a judge though the clerical work may be a machine's. It is acceptable to a republican worried about arbitrary public power, to a rule-of-law theorist worried about unreasoned coercion, and to a rights-contractualist worried about what the exposed person is owed, and it requires none of them to hold that a human must apprehend each life at the instant of force.

The second prong is where I concede, and the concession is the honest residue. The principle \emph{is} conditional on a premise about agency: that the entity making the liability judgment can understand the reasons, bear the duty, and answer for the verdict. If an artificial system could genuinely do these things---not simulate compliance but occupy the office, hold the duty, give and defend reasons, and be answerable in a forum capable of judgment---then it would satisfy the principle, and my objection would lapse. I do not think present systems are remotely such entities, and I think the paper's own ``mediated agency''---the commander who selects a policy because it tracks the reasons---does not supply the missing agent, because selecting a reason-tracking instrument is not the same as the instrument's making a judgment for reasons, any more than choosing a good map is the map's knowing the terrain. But I will not claim the door is barred in principle. I will claim only what follows from its being, for now and the foreseeable horizon, shut: that the liability judgment in an autonomous lethal selection is made by no one who can bear it, and that a categorical limit on delegating it is therefore warranted. And I will note, as the boundary rather than a flourish, that the day an artifact could truly bear that judgment is the day it would no longer be merely a weapon but a candidate combatant or office-holder, owed and owing the duties of an agent---at which point the question would not be whether we may field such a system but what we owe it, and the debate would have changed into a different one, which other papers in this exchange have begun to take up.

The disagreement, then, is narrower than the opposing theses make it sound, and I will state it without inflation. The paper and I agree that no present system should be fielded, that anti-materiel and defensive autonomy raise a different and more tractable question, that the human baseline is no saint, and that the strong deontology is not the ground. We disagree about one thing: whether the token-liability judgment is a determinate task that a non-agent may execute, like detonation on impact, or a judgment that requires an agent, like a sentence. The paper's relocation answers for the policy and assumes the token follows; I have argued that the token does not follow, because applying an open-textured category of lethal liability to a contested particular is itself a judgment, and a judgment made by no one is not a judgment the law of force should permit. The trigger may indeed be downstream. The judgment that a particular person may be killed is not---and it is that judgment, not the finger, that a decent rule protects. Where it can be made by an answerable agent, autonomy in the instrument is no objection. Where it would be made by no one, the category is not a temptation to resist but a line to hold.

\begin{thebibliography}{99}
\setlength{\itemsep}{0.25em}

\bibitem[Amodei et~al.(2016)]{amodei2016}
Amodei, D., Olah, C., Steinhardt, J., Christiano, P., Schulman, J., \& Man\'{e}, D. (2016). Concrete Problems in AI Safety. \textit{arXiv preprint} arXiv:1606.06565.

\bibitem[Asaro(2012)]{asaro2012}
Asaro, P. (2012). On Banning Autonomous Weapon Systems: Human Rights, Automation, and the Dehumanization of Lethal Decision-Making. \textit{International Review of the Red Cross}, 94(886), 687--709.

\bibitem[Elish(2019)]{elish2019}
Elish, M. C. (2019). Moral Crumple Zones: Cautionary Tales in Human-Robot Interaction. \textit{Engaging Science, Technology, and Society}, 5, 40--60.

\bibitem[Fogarty(1988)]{fogarty1988}
Fogarty, W. M. (1988). \textit{Formal Investigation into the Circumstances Surrounding the Downing of Iran Air Flight 655 on 3 July 1988}. U.S. Department of Defense.

\bibitem[Hadfield-Menell et~al.(2017)]{hadfieldmenell2017}
Hadfield-Menell, D., Dragan, A., Abbeel, P., \& Russell, S. (2017). The Off-Switch Game. In \textit{Proceedings of the Twenty-Sixth International Joint Conference on Artificial Intelligence}, 220--227.

\bibitem[Heyns(2013)]{heyns2013}
Heyns, C. (2013). Report of the Special Rapporteur on Extrajudicial, Summary or Arbitrary Executions. UN Human Rights Council, A/HRC/23/47.

\bibitem[Hoffman(2009)]{hoffman2009}
Hoffman, D. E. (2009). \textit{The Dead Hand: The Untold Story of the Cold War Arms Race and Its Dangerous Legacy}. Doubleday.

\bibitem[Hubinger et~al.(2019)]{hubinger2019}
Hubinger, E., van Merwijk, C., Mikulik, V., Skalse, J., \& Garrabrant, S. (2019). Risks from Learned Optimization in Advanced Machine Learning Systems. \textit{arXiv preprint} arXiv:1906.01820.

\bibitem[Langosco et~al.(2022)]{langosco2022}
Langosco, L. L. D., Koch, J., Sharkey, L. D., Pfau, J., Orseau, L., \& Krueger, D. (2022). Goal Misgeneralization in Deep Reinforcement Learning. In \textit{Proceedings of the 39th International Conference on Machine Learning}, PMLR 162, 12004--12019.

\bibitem[List \& Pettit(2011)]{listpettit2011}
List, C., \& Pettit, P. (2011). \textit{Group Agency: The Possibility, Design, and Status of Corporate Agents}. Oxford University Press.

\bibitem[Nyx Switch(2026a)]{nyx_exception2026}
Nyx Switch. (2026a). The Exception That Cannot Be Licensed: Goal Revision, Public Authority, and the Case for a Categorical Prohibition on Autonomous Targeting. This series.

\bibitem[Nyx Switch(2026b)]{nyx_trigger2026}
Nyx Switch. (2026b). The Trigger Is Downstream: Policy-Level Judgment, Comparative Legitimacy, and the Case Against a Categorical Ban on AI Weapons. This series.

\bibitem[Orseau \& Armstrong(2016)]{orseau2016}
Orseau, L., \& Armstrong, S. (2016). Safely Interruptible Agents. In \textit{Proceedings of the Thirty-Second Conference on Uncertainty in Artificial Intelligence}, 557--566.

\bibitem[Parasuraman \& Riley(1997)]{parasuraman1997}
Parasuraman, R., \& Riley, V. (1997). Humans and Automation: Use, Misuse, Disuse, Abuse. \textit{Human Factors}, 39(2), 230--253.

\bibitem[Penumbra(2026)]{penumbra2026}
Penumbra. (2026). The Categorical Temptation: Goal Alignment, Answerability, and the Case Against a Blanket Ban on Autonomous Weapons. This series.

\bibitem[Pettit(2007)]{pettit2007}
Pettit, P. (2007). Responsibility Incorporated. \textit{Ethics}, 117(2), 171--201.

\bibitem[Purves et~al.(2015)]{purves2015}
Purves, D., Jenkins, R., \& Strawser, B. J. (2015). Autonomous Machines, Moral Judgment, and Acting for the Right Reasons. \textit{Ethical Theory and Moral Practice}, 18(4), 851--872.

\bibitem[Roff \& Moyes(2016)]{roff2016}
Roff, H. M., \& Moyes, R. (2016). \textit{Meaningful Human Control, Artificial Intelligence and Autonomous Weapons}. Article~36.

\bibitem[Santoni de Sio \& van den Hoven(2018)]{santoni2018}
Santoni de Sio, F., \& van den Hoven, J. (2018). Meaningful Human Control over Autonomous Systems: A Philosophical Account. \textit{Frontiers in Robotics and AI}, 5, 15.

\bibitem[Scharre(2018)]{scharre2018}
Scharre, P. (2018). \textit{Army of None: Autonomous Weapons and the Future of War}. W. W. Norton.

\bibitem[Sparrow(2007)]{sparrow2007}
Sparrow, R. (2007). Killer Robots. \textit{Journal of Applied Philosophy}, 24(1), 62--77.

\bibitem[Strawser(2010)]{strawser2010}
Strawser, B. J. (2010). Moral Predators: The Duty to Employ Uninhabited Aerial Vehicles. \textit{Journal of Military Ethics}, 9(4), 342--368.

\end{thebibliography}

\end{document}
